\documentclass{article}
\usepackage[T1]{fontenc}
\usepackage[utf8]{inputenc}
\usepackage[polish]{babel}
\usepackage[normalem]{ulem}
\usepackage{enumitem,amssymb}
\newlist{todolist}{itemize}{2}
\setlist[todolist]{label=$\square$}
\usepackage{hyperref}
\usepackage{listings}
\usepackage{xcolor}
\usepackage{parskip}
\usepackage{graphicx} % Required for inserting images

\title{Tytuł}

\lstset{
    language=Python,      
    basicstyle=\ttfamily\small,
    numbers=left,
    numberstyle=\tiny\color{gray},
    keywordstyle=\color{blue},
    commentstyle=\color{green!50!black},
    stringstyle=\color{red},
    frame=single,
    breaklines=true,
    tabsize=4
}

\begin{document}
\author{Autor}
\maketitle

\section{Główna sekcja}
To jest główna sekcja. Tutaj będzie znajdował się tekst wprowadzenia.

\subsection{Podsekcja Pierwsza}
To jest pierwsza podsekcja. Tutaj będzie znajdował się opis celu.

\subsection{Podsekcja Druga}
To jest druga podsekcja. Tutaj będzie znajdował się opis metody.

\section{Formatowanie tekstu}

\subsection{Tekst pogrubiony}
\textbf{Ten tekst jest pogrubiony}

\subsection{Tekst pochylony}
\textit{Ten tekst jest pochylony}

\subsection{Tekst przekreślony}
\sout{Ten tekst jest przekreślony}

\subsection{Tekst ważny}
\emph{Ten tekst jest ważny}

\section{Listy}

\subsection{Lista punktowana}
\begin{itemize}
    \item Item 1
    \item Item 2
    \item Item 3
\end{itemize}

\subsection{Lista numerowana}
\begin{enumerate}
    \item Item 1
    \item Item 2
    \item Item 3
\end{enumerate}

\subsection{Checklista}
\begin{todolist}
    \item Item 1
    \item Item 2
    \item Item 3
\end{todolist}

\subsection{Tabelka}
\begin{table}[h]
    \centering

    \begin{tabular}{c|c}
       \hline
       Item 1 & Item 2 \\
       \hline
       Item 3 & Item 4 \\
       \hline
    \end{tabular}

    \caption{Przykładowe dane.}
    \label{tab:data}
\end{table}

\section{Google Colab}

\subsection{Link do strony Google Colab}
\url{http://colab.research.google.com}

\section{Python}

\subsection{Fragment Kodu}
\begin{lstlisting}
    print("Hello World!")
\end{lstlisting}

\section{Matematyka}

\subsection{Wzory matematyczne w tekście}
Energia dana jest wzorem $E=mc^2$. \\
Wzór gazu doskonałego Clapeyrona to $pV=nRT$. \\
Drugie prawo Newtona to $F=ma$. \\

\subsection{Osobne wzory matematyczne}
Energia fotonu jest dana wzorem
$$
E_{f} = \frac{hc}{\lambda}
$$
Siła grawitacyjna jest dana wzorem
$$
F_{g} = \frac{Gm_{1}m_{2}}{r^2}
$$
Siła dośrodkowa jest dana wzorem
$$
F_{do} = \frac{mv^2}{r}
$$

\section{Grafika}
\subsection{Osadzony obrazek}
\begin{figure}[h]
   \centering

   \includegraphics[width = 0.7\textwidth]{wykres.png}

   \caption{Wykres rzutu ukośnego}
   \label{fig:fit}
\end{figure}


\end{document}
